\let\negmedspace\undefined
\let\negthickspace\undefined
\documentclass[journal]{IEEEtran}
\usepackage[a5paper, margin=10mm, onecolumn]{geometry}
%\usepackage{lmodern} % Ensure lmodern is loaded for pdflatex
\usepackage{tfrupee} % Include tfrupee package

\setlength{\headheight}{1cm} % Set the height of the header box
\setlength{\headsep}{0mm}     % Set the distance between the header box and the top of the text

\usepackage{gvv-book}
\usepackage{gvv}
\usepackage{cite}
\usepackage{amsmath,amssymb,amsfonts,amsthm}
\usepackage{algorithmic}
\usepackage{graphicx}
\usepackage{textcomp}
\usepackage{xcolor}
\usepackage{txfonts}
\usepackage{listings}
\usepackage{enumitem}
\usepackage{mathtools}
\usepackage{gensymb}
\usepackage{comment}
\usepackage[breaklinks=true]{hyperref}
\usepackage{tkz-euclide} 
\usepackage{listings}
% \usepackage{gvv}                                        
\def\inputGnumericTable{}                                 
\usepackage[latin1]{inputenc}                                
\usepackage{color}                                            
\usepackage{array}                                            
\usepackage{longtable}                                       
\usepackage{calc}                                             
\usepackage{multirow}                                         
\usepackage{hhline}                                           
\usepackage{ifthen}                                           
\usepackage{lscape}
\begin{document}

\bibliographystyle{IEEEtran}
\vspace{3cm}

\title{12.151}
\author{EE25BTECH11033 - Kavin}
% \maketitle
% \newpage
% \bigskip
{\let\newpage\relax\maketitle}

\renewcommand{\thefigure}{\theenumi}
\renewcommand{\thetable}{\theenumi}
\setlength{\intextsep}{10pt} % Space between text and floats
\textbf{Question}:\\
A point $\vec{P}(10,3)$ MPa on the Mohr's circle represents normal and shear stresses. If the centre of the Mohr's circle is $\vec{C}(6,0)$ MPa, the normal and shear stresses in MPa on the point diametrically opposite to $\vec{P}$ are

\begin{enumerate}
	\begin{multicols}{4}
    \item 2, $-3$
    \item 4, $-3$
    \item 2, 3
    \item 4, 3
	    \end{multicols}
\end{enumerate}\\    
\bigskip

\textbf{Solution}:\\
A point on the Mohr's circle, 
\begin{align}
    \vec{P} = \myvec{\sigma_P\\ \tau_P} = \myvec{10\\3} \text{ MPa}
\end{align}
where,
\begin{align*}
    \sigma = \text{Normal Stress} \text{ and } \tau = \text{Shear Stress}
\end{align*}
The center of the Mohr's circle,
\begin{align}
    \vec{C} = \myvec{\sigma_C\\ \tau_C} = \myvec{6\\0} \text{ MPa}
\end{align}
Let P' be the point diametrically opposite to P ,
\begin{align}
    \vec{P'} = \myvec{\sigma_P'\\ \tau_P'} \text{ MPa}
\end{align}
The center $\vec{C}$ is the midpoint of the line segment connecting $\vec{P}$ and $\vec{P'}$. We can use the midpoint formula for both the horizontal (normal stress) and vertical (shear stress) coordinates.
\begin{align}
    \vec{C} = \frac{\vec{P} + \vec{P'}}{2}
\end{align}
\begin{align}
    \vec{P'} = 2\vec{C}-\vec{P}
\end{align}
\begin{align}
    \vec{P'} = 2 \myvec{6\\0}-\myvec{10\\3}
\end{align}
\begin{align}
    \implies \vec{P'} = \myvec{\sigma_P'\\ \tau_P'} = \myvec{2\\-3}\text{ MPa}
\end{align}

\newpage
\begin{figure}[H]
\begin{center}
\includegraphics[width=0.7\columnwidth]{figs/fig.png}
\end{center}
\label{fig:Fig1}
\end{figure}
\end{document}


